\documentclass[11pt,a4paper]{article}
\usepackage[margin=2cm]{geometry}
\usepackage[T1]{fontenc}
\usepackage{mathptmx}
\usepackage{amsmath}
\usepackage{xcolor}
\usepackage{enumitem}
\usepackage{array}
\usepackage[hidelinks]{hyperref}
\definecolor{navy}{HTML}{1B3755}
\definecolor{blue}{HTML}{2F6FB3}
\definecolor{red}{HTML}{C62828}
\setlength{\parindent}{0pt}
\setlength{\parskip}{4pt}
\newcounter{qn}
% \Q{question}{answer}
\newcommand{\Q}[2]{\refstepcounter{qn}\par\medskip
  {\color{navy}\bfseries Q\theqn.\ #1}\par\nopagebreak
  \hspace*{1.2em}\begin{minipage}[t]{\dimexpr\linewidth-1.2em}#2\end{minipage}\par}
% tricky question with a "trap" note
\newcommand{\T}[3]{\refstepcounter{qn}\par\medskip
  {\color{red}\bfseries Q\theqn.\ #1}\par\nopagebreak
  \hspace*{1.2em}\begin{minipage}[t]{\dimexpr\linewidth-1.2em}{\color{red}\itshape Why it is tricky: #2}\par #3\end{minipage}\par}
\newcommand{\sect}[1]{\bigskip{\Large\color{navy}\bfseries #1}\par\smallskip{\color{navy}\hrule height 0.8pt}\smallskip}

\begin{document}
\begin{center}
{\LARGE\bfseries\color{navy} Presentation Defence Preparation}\\[4pt]
{\large Dual-Setting Directional Recloser and Fuse Coordination with DG --- IEEE 13-Node Feeder, DIgSILENT PowerFactory}\\[2pt]
Jhala Nath Kafle (081MSPSE009)
\end{center}

\sect{1. Numbers to Remember}
\begin{tabular}{|>{\raggedright}p{6.3cm}|>{\raggedright\arraybackslash}p{9.6cm}|}
\hline
Feeder & IEEE 13-node, 4.16\,kV, unbalanced; 115\,kV grid, 5\,MVA substation transformer \\ \hline
DG & 4.05\,MVA synchronous machine at node 692, 0.69/4.16\,kV transformer ($u_k$ 15\,\%), 3.24\,MW \\ \hline
R1 (feeder head, RG60--632) & GE IAC77B801A, extremely inverse, pickup 720\,A, TDS 0.5 (fast) / 10 (delayed) \\ \hline
R2 forward (line 632--671) & CDG34, plugs 300 / 600\,A, TMS 0.1 / 1.0 \\ \hline
R2 reverse & CT 500/5, plugs 150 / 300\,A, TMS 0.1 / 1.0; reverse pickup $1.25\times252$\,A $\approx 315$\,A \\ \hline
Study grid & 4 fault types at 12 locations = 39 cells \\ \hline
Coordination held & no DG 39 (35 before F692-R $\rightarrow$ 200E); DG + conventional R2 24; DG + DSDR + revised fuses 38; same fuses, single setting 31 \\ \hline
Fuse revision & F632 500E; F633, F646, F-DL, F671-1 400E; F692-R 250E \\ \hline
Still lost & LG fault at 692 (dial limit of R2) \\ \hline
Fault level check & IEEE benchmark within about 2\,\%; feeder head 4.73\,kA \\ \hline
PowerFactory relay check & 1896 operating times, largest deviation 1.22\,\% \\ \hline
EMT (LL a--c at 684, 0.2\,$\Omega$) & no DG: fuse melts 1.248\,s, clears 1.613\,s (fast shots use 45\,\% of melting heat); with DG: melts 0.752\,s, clears 1.253\,s \\ \hline
DG penetration (fixed settings) & CTI 633: $+4 \rightarrow -51$\,ms; 671: $+403 \rightarrow +75$\,ms; fault level up to +64\,\% \\ \hline
Margins & breaker time 3 cycles: 31 cells; plus fuse safety margin: 26 cells \\ \hline
\end{tabular}

\sect{2. Limitations (and how to present them)}
Say the limitation first, then why it does not invalidate the result. Never hide one; a panel trusts a
candidate who names the weak points before being asked.
\begin{enumerate}[leftmargin=*,itemsep=3pt]
\item \textbf{Unpublished data.} The paper does not give R2's settings, the exact fuse sizes or the source
  impedance. \emph{Answer:} they were chosen by the paper's own equations and every choice is documented in the
  report, so the work is reproducible.
\item \textbf{CDG34 is not directional in the PowerFactory library.} The two setting groups are two relay units,
  and the direction is assigned from the fault location (topology). \emph{Answer:} for a radial feeder with one DG
  the direction of the current through R2 is fixed by which side the fault is on, so this is exact for the
  study; a real relay would use a directional element (67), listed as future work.
\item \textbf{Zero-margin criterion.} The fuse must not start to melt before the fast trip. \emph{Answer:} this is
  the paper's rule; the sensitivity was checked: with a 3-cycle breaker time 31 cells, with a fuse margin 26.
\item \textbf{Dial ranges are fixed by the real devices.} The IAC and CDG curves cannot be made faster or slower
  beyond their ranges, which is why one cell (LG at 692) stays lost. \emph{Answer:} this is a property of real
  hardware that an idealised curve would hide --- a finding, not an error.
\item \textbf{One DG location and one DG type.} Only a synchronous DG at 692. Inverter-based DG gives much
  less fault current (about 1.1--1.5 p.u.) and is future work.
\item \textbf{Faults downstream of R2.} The DG feeds them directly with no recloser in its path, so no setting of
  R2 can save those fuses. \emph{Answer:} this is a limit of the DSDR concept itself, shown by the time-sequence check.
\item \textbf{Fixed settings in the penetration study.} The settings of the design without DG were kept at every
  level to show the degradation; an adaptive scheme would re-set at each level.
\item \textbf{Fault levels lower than the paper.} The model was validated against the IEEE benchmark instead
  of the paper's Table II (whose values are internally inconsistent, see Q26).
\item \textbf{IEEE 34-node feeder not modelled} (in the paper it is a second test system).
\item \textbf{Single EMT case.} Only one time-domain case (LL at 684); the rest are curve-based calculations,
  confirmed by PowerFactory's own relay models within 1.22\,\%.
\end{enumerate}

\sect{3. Basic Questions}
\Q{What is fuse saving?}{The recloser trips on its fast curve before the lateral fuse melts, recloses, and a
temporary fault (most faults on overhead lines) is cleared without blowing the fuse. If the fault is permanent,
the delayed curve gives the fuse time to clear, so only the faulted lateral is lost.}

\Q{What exactly does the DG change?}{Two things. (1) Magnitude: the fuse carries grid current plus DG current,
the recloser only the grid current, so the fuse becomes relatively faster. (2) Direction: for a fault between the
grid and the DG, the DG current flows backwards through R2 (from 671 to 632).}

\Q{What is a DSDR?}{A recloser with two independent setting groups, forward and reverse, each with its own pickup
and time dial; the group is chosen by the direction of the current. R2 can then be fast enough for the small
reverse DG current and still coordinate with the fuses downstream for forward faults.}

\Q{Why is R2 placed on line 632--671?}{It separates the grid side (632, 633, 645, 646) from the DG side (671,
692, 675, 680, 684). It is the only recloser that sees current in both directions.}

\Q{Why IEEE 13-node feeder?}{Short, highly loaded, unbalanced, with single-, two- and three-phase laterals,
a regulator and a transformer --- a hard test for protection, and the system used in the reference paper.}

\Q{How did you calculate the pickup current?}{$I_p = \mathrm{OLF}\times I_{nom}$ with OLF = 1.25 and $I_{nom}$ the
load current from the load flow. The pickup must also be below the smallest fault current of the zone (LG through
3\,$\Omega$ at the farthest node) --- 1074\,A for R1 and 879\,A for R2 without DG.}

\Q{Why is the reverse pickup lower than the forward one?}{The reverse current is only the DG's contribution, much
smaller than the grid's. The reverse load current through R2 with the DG is 252\,A, so $1.25\times252 \approx
315$\,A. A forward setting of 300/600\,A plugs would be too slow or would not even start for this current.}

\Q{What does ``cell'' mean in your results?}{One combination of fault location and fault type. Four types at twelve
locations, minus combinations that cannot exist on single- and two-phase laterals, gives 39 cells.}

\Q{What is the CTI and why does it matter?}{The coordination time interval, $t_{MMT}(\text{fuse}) - t_{fast}
(\text{recloser})$. Positive: the recloser trips before the fuse starts to melt (fuse saved). Negative: the fuse
melts first (fuse saving lost).}

\Q{What is the series-fuse condition, eq.\ (8)?}{When two fuses are in series, the fuse nearer the fault must clear
before the upstream fuse starts to melt; the usual rule is: total clearing time of the downstream fuse $\le 0.75
\times$ minimum melting time of the upstream fuse. The 0.75 leaves room for preheating and tolerances.}

\Q{Why DIgSILENT PowerFactory?}{It has the library curves of the real relays (IAC77, CDG34) and fuses (A055C),
unbalanced load flow and short circuit, and EMT simulation in one tool, and its own relay models were used to
check the calculated times (1896 times, within 1.22\,\%).}

\sect{4. Method and Result Questions}
\Q{Walk me through the method.}{Stage A: conventional design without DG (pickups, dials, fuse sizes) --- 39 of 39
after one fuse change. Stage B: same settings, DG connected --- 24 of 39. Stage C: R2 as DSDR; dials could not
help (both relays at their limits), so the fuses were revised --- 38 of 39.}

\Q{How much of the improvement is the DSDR and how much the fuse revision?}{Separated on purpose: DSDR with the old
fuses 25 (+1); revised fuses with a single-setting R2 31 (+7); both together 38 (+14). So seven of the fourteen
restored cells need the dual setting; the two measures work together.}

\Q{Which cells does the dual setting restore?}{All four fault types at 633 and the LG faults at 645, 646 and the
distributed load --- all in the zone between the grid and the DG, where R2 sees reverse current.}

\Q{Why is the LG fault at 692 still lost?}{F671-1 must be 400E to stay selective with F692-R, and then for the LG
fault at 692 the fast curve of R2 would have to be faster than its smallest dial (TMS 0.1) allows. It is a limit
of the real device's dial range.}

\Q{Explain your single vs dual example (three-phase fault near 632).}{R2 carries 1777\,A in reverse (DG only). Single
setting: forward curve, trips at 0.121\,s, but F633 starts to melt at 0.100\,s --- lost. Dual setting: reverse group,
trips at 0.052\,s --- held. Same fault, same fuse; only the setting group changed.}

\Q{What does the time-sequence check add?}{The classification uses a strict instant rule. The time-sequence check
follows each fault in time and lets the fuse heat accumulate while current flows after each recloser opens. Above
R2: 15 of 18 held with a single setting, 18 of 18 with the DSDR. Below R2 the DG keeps feeding the fault while R2
is open, so the setting of R2 does not matter there.}

\Q{What does the EMT simulation show?}{LL (a--c) fault at 684 through 0.2\,$\Omega$, conventional settings. Without DG
the two fast shots use 45\,\% of the fuse's melting heat; the fuse survives and clears the permanent fault at
1.61\,s. With DG the DG keeps feeding the fault during the dead time; the fuse melts at 0.75\,s, during the second fast
shot --- fuse saving lost.}

\Q{In the EMT figure the fault is at 684. Why is the current plotted at node 632?}{Node 632 is on the path of the
fault current (substation, 632, line 632--671 with R2, 671, F671-2, 684) and shows the whole sequence: each fast shot
of R2, the reclosures, the fuse clearing, and the healthy feeder supplied again afterwards. At 684 the current simply
stops when the fuse opens, so the restoration would not be visible. The paper uses the same point.}

\Q{Why a fault at 684 and not on the line 671--684?}{PowerFactory refused the line-fault event on that two-phase
section, so the fault was applied at node 684, at the end of the same line; the fuse F671-2 and R2 see the same
path.}

\Q{Line 671--684 is two-phase (a, c). Why is the fault called ``a--b'' in PowerFactory?}{PowerFactory names the phases
of a two-phase element locally (first and second conductor). On this line the first is IEEE phase a and the second
IEEE phase c, so ``a--b'' in PowerFactory is the a--c fault of the study.}

\Q{Why does the CTI fall as DG penetration rises?}{The fuse is between the fault and both sources, so it carries
grid plus DG current; the recloser carries only the grid share, which even falls slightly because the DG holds the
voltage up. The fuse curve is steep, so its melting time drops much faster than the recloser time rises. 633: $+4
\rightarrow -51$\,ms; 671: $+403 \rightarrow +75$\,ms.}

\Q{How did you verify that your calculated operating times are right?}{The settings were entered into PowerFactory's
relay and fuse elements and PowerFactory calculated the tripping time of every device for every fault: 1896 times,
largest deviation 1.22\,\%. Also the R1 worked example of the paper is reproduced (0.096 / 1.926\,s vs 0.097 /
1.932\,s).}

\sect{5. Tricky Questions}
\T{Your DSDR is not really directional --- the library relay has no directional element. Is this a DSDR at all?}
{It attacks the core claim.}
{The setting groups and their coordination are what the study tests; the direction decision is a separate, simple
function. In a radial feeder with one DG the direction of the current in R2 is determined by the side of the fault,
so assigning the group by fault location gives exactly what an ideal directional element would. The weakness is
stated in the limitations, and an explicit directional element is the first item of future work.}

\T{The paper gets 39 of 39 with the DSDR, you get 38. Did your method fail?}{Invites you to admit a mistake.}
{No. The paper does not publish R2's settings or all fuse sizes. With the real CDG34 dial range, the LG fault at 692
needs a faster curve than TMS 0.1 gives. That is a finding about real hardware. The central finding --- the DSDR
restores coordination between the grid and the DG --- is reproduced.}

\T{Your fault currents are lower than the paper's. Is your model wrong?}{Questions model validity.}
{The model was validated against the IEEE short-circuit benchmark and agrees within about 2\,\% (4.73\,kA at the head).
The paper's Table II cannot be fully right: it gives 6.73\,kA on 632--633 and 7.83\,kA on 692--675, both above its own
5.41\,kA at the feeder head, which is impossible for a radial feeder. The rated load currents agree within 2\,\%.}

\T{You changed the fuse sizes. Isn't that cheating --- any coordination can be fixed with bigger fuses?}
{Suggests the result is artificial.}
{Fuse upgrade is step 9 of the paper's own method, used only after the dial revision cannot help. And bigger fuses
alone are not enough: with the revised fuses and a single-setting R2 only 31 cells are held; the dual setting is
needed for the other seven. Bigger fuses also have a cost (slower clearing of impedance faults), which is why cases
3 and 4 of the nine-fault set are lost with both settings.}

\T{With a zero margin your 38 cells are optimistic. What happens with a realistic margin?}{Tests honesty.}
{Correct, and it was checked: adding a 3-cycle breaker time gives 31 cells, adding also a fuse safety margin gives 26.
The relative benefit of the DSDR remains; the absolute numbers depend on the margin, which the paper does not state.}

\T{Your progress report said the CTI increases with DG; now it decreases. Which is right?}{Exposes an earlier
inconsistency.}
{The decreasing trend is right and it is physically explained (fuse carries grid + DG, recloser only grid). The
progress-stage model was built in MATLAB with a static Thevenin DG and recalibrated fuse coefficients. The final model
in PowerFactory uses a dynamic synchronous machine and library fuse curves, and agrees with the paper's trend. The
change of tool was made exactly because the earlier result could not be explained.}

\T{The reverse pickup is 315\,A, but the reverse fast plug is 150\,A. Which one is the pickup?}{Looks like a
contradiction.}
{The CDG34 time curve is defined from a plug-setting multiplier of 2, so the curve starts at twice the plug: 150\,A
fast plug $\rightarrow$ 300\,A start, 300\,A delayed plug $\rightarrow$ 600\,A start. 300\,A is the closest available
setting to the calculated 315\,A pickup.}

\T{With more DG, couldn't the reverse load current exceed the reverse pickup and trip R2 on load?}{Probes a hidden
assumption.}
{Yes, that is why the method's last box says the settings are valid for the DG capacity they were designed for. At
4.05\,MVA the reverse load is 252\,A, below the pickup. A larger DG would need new reverse settings --- which is
exactly what an adaptive, setting-group scheme allows.}

\T{During the reclosing dead time, wouldn't the DG trip by anti-islanding protection anyway?}{Challenges the
EMT scenario.}
{Anti-islanding standards (e.g.\ IEEE 1547) allow up to 2\,s to cease energising; the dead time here is 0.2\,s, so the
DG is still connected during the fast reclosing shots --- that is the dangerous case shown. DG disconnection logic
during the dead time is listed as future work.}

\T{Why not just disconnect the DG for every fault instead of a DSDR?}{Asks whether the work is needed at all.}
{Disconnecting the DG for every temporary fault loses its generation and reliability benefit and needs fast transfer
trip communication. The DSDR uses only local measurements at R2 and keeps the DG in service.}

\T{Would the same work for solar inverters?}{Extrapolation beyond the study.}
{Not directly. Inverters limit fault current to about 1.1--1.5 times rated, so the fuse--recloser mismatch is
smaller, but detecting the direction becomes harder. That is future work; this study is for a synchronous DG.}

\T{The paper's fuse equation with its Table III doesn't give the times in its own figures. How do you know your
fuse curves are right?}{Tests whether you checked the source.}
{The library A055C curves were used, not the equation. For F646 at 3.6\,kA the paper's figure shows about 0.18\,s,
while its equation with Table III gives 1.85\,s; the figures match the A055C library curves. Table III could only be
reproduced with the series index reversed.}

\T{Your CTI at node 633 is negative already without much DG. Was the design bad from the start?}{Suggests a poor
base design.}
{Without DG the margin at 633 is only +4\,ms, so it turns negative from about 4\,\% penetration. That small margin
comes from the fuse size of the design without DG. It shows how sensitive a tight design is --- the reason the method
revises the fuses (F633 to 400E).}

\T{What is new in your work compared with the paper?}{Asks for your contribution.}
{(1) An independent implementation in PowerFactory with real relay and fuse library curves and a dynamic DG.
(2) Separation of the DSDR effect from the fuse-revision effect (25 / 31 / 38 cells). (3) A time-sequence check of
all cells and an EMT check. (4) A margin sensitivity study. (5) Verification of every operating time with
PowerFactory's own relay models. (6) A documented check of the paper's tables.}

\sect{6. Tips for the Defence}
\begin{itemize}[leftmargin=*,itemsep=2pt]
\item Answer with a number from your results whenever possible (24, 31, 38; 0.121 vs 0.052\,s).
\item If you do not know: ``I did not study that; I would expect \dots\ because \dots'' --- never guess a number.
\item When a weakness is raised, agree, then give the check you made (margin study, PowerFactory relay check).
\item Keep the comparison with the paper for the appendix slide unless the panel asks earlier.
\item Have the report open at the results summary table and the comparison table.
\end{itemize}
\end{document}
